\section{Open Challenges}
\label{sec:challenges}
\begin{table*}[t]
\centering
\caption{Open interfaces between navigation functions, the dependency that each leaves unresolved, and a matched comparison that would test it. Each interface lists its coded works and those with direct joint wheeled evidence.}
\label{tab:roadmap}
\surveytableformat
\begin{tabularx}{\linewidth}{@{}L{2.3cm}L{3.6cm}YL{4.0cm}@{}}
\toprule
\textbf{Interface} & \textbf{Unresolved dependency} & \textbf{Matched comparison} & \textbf{Primary outcomes} \\
\midrule
1. Localization-triggered recovery (Sec.~\ref{sec:challenges}-A). \IfaceOneWorks{} works, \IfaceOneDirect{} direct & Which response a localization warning should select when part of the state needed to execute it is unreliable & Fixed estimator and route. Immediate stop, threshold-triggered relocalization, and an action chosen from the remaining trustworthy channels, under controlled pose bias and weak geometry, including trials in which one alarm fires while another channel misleads & Restored localization and task progress, contact and interventions, cost of unnecessary stops under nominal sensing \\
\addlinespace[3pt]
2. Semantic reasoning with fallback (Sec.~\ref{sec:challenges}-B). \IfaceTwoWorks{} works, \IfaceTwoDirect{} direct & How long a semantic decision stays valid, and whether wrong as well as late responses are rejected & Fixed model and controller. Periodic querying, scene-triggered querying, and cached decisions, under delayed responses, changed candidate paths, injected incorrect responses, and varied obstacle-track quality & Completion, contact, and fallback duration against energy and query cost, with assignment age for teams \\
\addlinespace[3pt]
3. Pose uncertainty in terrain decisions (Sec.~\ref{sec:challenges}-C). \IfaceThreeWorks{} works, \IfaceThreeDirect{} direct & How pose uncertainty changes the probability that the wheel footprint rests on traversable support & Fixed terrain map and chassis. A calibrated pose distribution propagated into footprint clearance or speed, against geometry-only planning and fixed inflation, on narrow supports and slopes with an independent reference & Contact and interventions against completion and travel time, with aborted risky attempts counted as outcomes \\
\addlinespace[3pt]
4. Shared-state acceptance after reconnection (Sec.~\ref{sec:challenges}-D). \IfaceFourWorks{} works, \IfaceFourDirect{} direct & Which returning information to accept when it may be stale, correlated, or biased & Moving ground team after a controlled partition. Immediate fusion, age-based rejection, and correlation- or residual-aware acceptance against a deterministic merge, with delay and landmark bias varied separately & Fault propagation and shared-state divergence, duplicated coverage, completion, and interventions \\
\bottomrule
\end{tabularx}
\end{table*}
The four interfaces below, numbered as in Fig.~\ref{fig:stack} and specified with a matched comparison in Table~\ref{tab:roadmap}, connect these functions in ways whose behavior under degradation remains insufficiently characterized. In the coded evidence, every pair of functions with one coupling source forms a candidate, ranked by a pre-registered rule on the number of its works without direct joint wheeled evidence, that is, without a named condition of any source and a physical task outcome in one experiment on a target wheeled robot (Section~\ref{S-supp:interfaces} of the supplementary material). Among the \CandidateCount{} candidates with at least three works, Interfaces~1--4 rank \IfaceOneRank{}, \IfaceTwoRank{}, \IfaceThreeRank{}, and \IfaceFourRank{}. Without the integration-questions stage, which was acquired for them, they rank \IfaceOneRankA{}, \IfaceTwoRankA{}, \IfaceThreeRankA{}, and \IfaceFourRankA{}. The interfaces therefore remain an agenda proposed by this survey rather than the four highest-ranked candidates. Six of the \RejectedCount{} candidates that rank above Interface~3 without being developed are merged into Interface~4 (team communication, shared-state faults, and teammate perception) or Interface~2 (moving agents in semantic choices and learned decisions). Section~\ref{sec:challenges}-E names the seventh.

\subsection{Localization-Triggered Navigation Recovery}
Localization monitoring and executed recovery are both available, although the reviewed systems rarely connect them through the same trigger. Monitors cover pose-error protection levels~\cite{P_F01}, modality switching~\cite{P_F04}, landmark repair on road-vehicle data~\cite{P_F03}, and fault-tolerant tracking of specified channels~\cite{P_F02}. Executed responses include stop--reposition--resume after optimizer failure~\cite{P_E59}, stopping and relocalization after visual tracking loss~\cite{P_N10}, behavior-tree recovery on CENTAURO~\cite{P_N17}, view planning against tracking loss~\cite{P_R33}, and interoceptive wheel--inertial odometry whose drift grows without external correction~\cite{P_R04}. Among the \IfaceOneWorks{} coded works of this interface, only the cross-modal teach-and-repeat study gives direct joint wheeled evidence~\cite{P_R11}, whereas the integrity analyses rest on road-vehicle data~\cite{P_X07,P_M28}. These triggers concern different objects, namely optimizer infeasibility, predicted tracking loss, and hazardous pose error. The open question is how a localization warning should select slowing, stopping, relocalization, or repositioning when part of the state needed to execute that action is itself unreliable. In automated driving, the corresponding response is a minimum risk maneuver toward a minimum risk condition~\cite{B_MRMFallback2025}, whereas a wheeled robot can often also relocalize or reposition within its drivetrain limits (Table~\ref{tab:drivetrain}).

\subsection{Resource-Constrained Semantic Reasoning and Fallback}
Wheeled systems already combine semantic trajectory selection with geometric planning~\cite{P_F05,P_R25}, group reasoning with constrained local control~\cite{P_F06}, and embodiment-specific re-ranking~\cite{P_R26}, while perception calibration on a legged robot~\cite{P_F08}, runtime monitoring~\cite{P_F07}, cached reasoning with fallback on a quadrotor~\cite{P_E60}, and distilled encoders~\cite{P_R29} address reliability and response time. Language-derived costmap layers~\cite{P_R30} and relocation-aware object search~\cite{P_R27} move semantic errors into costs or search order, where geometry can check them. This interface has the most direct wheeled evidence of the four, \IfaceTwoDirect{} of its \IfaceTwoWorks{} coded works, although none of them names a computation limit or response delay among its conditions. The remaining interface concerns the lifetime of a semantic decision, since an old destination can remain valid while an old pedestrian detour becomes unusable. A delayed team assignment, as in Co-NavGPT with perception and planning on a workstation~\cite{P_N09}, can likewise remain syntactically valid after the shared map or the allocation has changed. HAM-Nav and Mobility VLA depend on sketch or tour coverage, while VLM-GroNav retains GPS and cloud dependencies~\cite{P_S39,P_S43,P_S40}. SocialNav-SUB indicates that incorrect reasoning, and not only unavailable or delayed responses, requires explicit testing~\cite{P_S44}.

\subsection{Pose Uncertainty in Wheeled Terrain Decisions}
Recent terrain decisions increasingly combine geometry with localization, map, or contact uncertainty, through localizability-aware traversability~\cite{P_E13}, terrain-aware Ackermann allocation~\cite{P_F09}, uncertainty-weighted mapping on an aquatic explorer~\cite{P_F10}, uncertain elevation and terrain classes~\cite{P_S46,P_R08}, novelty and speed distributions~\cite{P_S42}, and angle-dependent risk~\cite{P_S45}. Terrain-adaptive stability bounds~\cite{P_R13} and implicit traversability fields~\cite{P_R15} shape executed speed and route on small skid-steer robots, whose hardware failures are attributed to wheel slip, wet-ground traction, and sensor occlusion rather than to pose error. STEP comes closest to this question among these planners, since localization error enters its collision risk through the elevation-map variance, although its simulation assumes no localization error and its field runs on a Husky and a Spot quadruped report no failure rates~\cite{P_K03}. EVORA instead models traction uncertainty under motion-capture pose indoors~\cite{P_K05}. These terrain planners give direct wheeled evidence for terrain risk without coding localization as a function, which leaves only \IfaceThreeWorks{} works in the coded cells of this interface. Their single direct wheeled study, an OpenStreetMap-grounded navigator, is direct through injected map deviations rather than pose uncertainty~\cite{P_M10}, which leaves this question with extrapolated evidence only. The open question is how uncertainty in pose, beyond its share of map variance, alters the probability that a wheel-supported footprint lies on the estimated traversable surface under incomplete observation, and whether propagating it changes executed outcomes.

\subsection{Disconnection, Reconnection, and Corrupted Shared Information}
Reconnection mechanisms exist for particular shared quantities, such as target estimates fused after a link returns in failure-aware aerial tracking~\cite{P_F12}, cooperative pose fusion~\cite{P_F15}, and correlation-aware landmark repair, a candidate estimator-level mechanism rather than a validated reconnection protocol~\cite{P_F03}. Reception-aware control~\cite{P_F13}, sparse position sharing~\cite{P_F14}, and rendezvous co-optimized with allocation~\cite{P_R19} change connectivity or message requirements. Asynchronous consensus tolerates dropped exchanges and outliers~\cite{P_R20,P_R23}, whereas a deterministic regional merge tolerates delay without testing corrupted content~\cite{P_R22}. Neither outlier rejection of inter-robot loop closures~\cite{P_K06,P_K07} nor store-and-forward sharing in the Subterranean Challenge teams~\cite{P_K08,P_K09,P_K10} isolates the effect of stale or biased batches. Sparse loop exchange, lossy keyframe coding, layered transmission, and planned check-ins reduce traffic~\cite{P_S10,P_S29,P_R21,P_S25} without establishing safe acceptance after reconnection. None of the \IfaceFourWorks{} coded works of this interface gives direct joint wheeled evidence, which places it first in the ranking. The \IfaceFourWhl{} target wheeled works among them name conditions only in recorded data, replay, online estimation, or simulation, if at all, which leaves this interface on extrapolated evidence. The open interface is selective acceptance beyond outlier rejection, since receiving a packet establishes neither frame compatibility nor the validity of its map or pose content. Concrete prerequisites to violate one at a time include stable metric odometry in foundation-model loop fusion~\cite{P_S11}, continuous range visibility in cooperative initialization~\cite{P_S27}, orchard-specific overlap priors~\cite{P_S23}, and stable front-end odometry in asynchronous collaborative LiDAR SLAM~\cite{P_R20}.

\subsection{Broader Challenges}
Interacting degradation sources call for benchmarks with controlled severity and onset and repeated independent missions. Transfer across terrain, supervision, platforms, and simulators~\cite{P_M42,P_E04,P_R15,P_R36} calls for tests on withheld conditions. Computation and energy choices~\cite{P_M11,P_E60,P_R31,P_R32} require response-time distributions under concurrent loads. Map maintenance, long deployments, and continual learning~\cite{P_E16,P_R07,P_R11,P_E15,P_M37,P_S14,P_S41,P_R29} must be related to revisits, availability, and human escalation, which \HumanAssistCount{} coded works note as human assistance outside the ranked function pairs. Assurance methods~\cite{P_X07,P_F08,P_E25,P_R14} need online monitors of their own assumptions. Embodiment-specific affordances and failure-aware team policies~\cite{P_R26,P_R35} call for cross-platform tests of clearance, curvature, and traction. Among the couplings ranked above Interface~3, actuator faults accommodated inside one tracking controller~\cite{P_M15,P_E52,P_F02} still lack tests in which planning or supervision responds to them.
